\documentclass[aspectratio=169]{beamer}

\input{preamble.tex}

\title[Rust статический полиморфизм]
{Rust статический полиморфизм}

\institute[СПбГУ]{}

\author[Дмитрий Козенко]{Козенко Дмитрий Сергеевич, группа 23.Б08-мм}

\begin{document}

\begin{frame}[fragile]{Условная компиляция}
\begin{minted}[fontsize=\small]{rust}
fn greeting() -> &'static str {
  #[cfg(feature = "verbose")]
  {
    #[cfg(target_os = "linux")] { return "linux"; }
    #[cfg(target_os = "windows")] { return "windows"; }
    #[cfg(not(any(target_os = "linux", target_os = "windows")))]
    { return "unknown"; }
  }
  #[cfg(not(feature = "verbose"))] { return ""; }
}

fn main() {
  println!("hello {}", greeting());
}

// cargo run --features verbose # -> hello linux
// cargo run # -> hello
\end{minted}
\end{frame}

% ---------------------------------------------------------------------

\begin{frame}[fragile]{Классическое использование трейтов}
\begin{minted}[fontsize=\small]{rust}
trait Animal { fn speak(&self); }

struct Dog;
struct Cat;

impl Animal for Dog { fn speak(&self) { println!("Woof!"); } }
impl Animal for Cat { fn speak(&self) { println!("Meow!"); } }

// Wlll be generated 2 functions
fn announce<T: Animal>(animal: T) { animal.speak(); }

fn main() {
  announce(Dog); // Woof
  announce(Cat); // Meow
}
\end{minted}
\end{frame}

% ---------------------------------------------------------------------

\begin{frame}[fragile]{Реализация по умолчанию}
  \begin{columns}[T, onlytextwidth]
    \setlength{\columnsep}{0pt}
    \begin{column}{0.5\textwidth}
      \begin{minted}[fontsize=\footnotesize]{rust}
trait Animal { fn speak(&self); }

impl<T> Animal for T {
  fn speak(&self) { println!("Hello"); }
}

struct Cat;

fn announce<T: Animal>(animal: T) {
  animal.speak();
}

fn main() {
  announce(Cat); // Hello
}
      \end{minted}
    \end{column}
    \begin{column}{0.5\textwidth}
      \begin{minted}[fontsize=\footnotesize]{rust}
trait Animal { fn speak(&self); }
impl<T> Animal for T {
  fn speak(&self) { println!("Hello"); }
}
struct Cat;
impl Animal for Cat { fn speak(&self)
  { println!("Meow!"); } }
fn announce<T: Animal>(animal: T)
  { animal.speak(); }
fn main() {
  announce(Cat);
}
//  3 | impl<T> Animal for T {
//    | first implementation here
// 10 | impl Animal for Cat { fn speak(&self)
//    | conflicting implementation for `Cat`
      \end{minted}
    \end{column}
  \end{columns}
\end{frame}

% ---------------------------------------------------------------------

\begin{frame}[fragile]{feature(specialization) --- nightly версия}
  \begin{columns}[T]
    \begin{column}{0.48\textwidth}
      \begin{minted}[fontsize=\small]{rust}
// rustup override set nightly
#![feature(specialization)]
trait Animal {
  fn speak(&self);
}
impl<T> Animal for T {
  default fn speak(&self) {
    println!("Hello");
  }
}
struct Cat;
struct Dog;

      \end{minted}
    \end{column}
    \begin{column}{0.48\textwidth}
      \begin{minted}[fontsize=\small]{rust}

impl Animal for Cat {
  fn speak(&self) {
    println!("Meow!");
  }
}
fn announce<T: Animal>(animal: T) {
  animal.speak();
}
fn main() {
  announce(Cat); // Meow!
  announce(Dog); // Hello
}
      \end{minted}
    \end{column}
  \end{columns}
\end{frame}

% ---------------------------------------------------------------------

\begin{frame}[fragile]{Ассоциированные типы}
  \begin{columns}[T, onlytextwidth]
    \setlength{\columnsep}{0pt}
    \begin{column}{0.5\textwidth}
      \begin{minted}[fontsize=\footnotesize]{rust}
trait Container {
    type Item;
    fn get(&self) -> Option<&Self::Item>;
}

struct VecWrapper<T> {
  data: Vec<T>,
}
impl<T> Container for VecWrapper<T> {
  type Item = T;
  fn get(&self) -> Option<&Self::Item> {
    self.data.first()
  }
}

fn main() {}
      \end{minted}
    \end{column}
    \begin{column}{0.5\textwidth}
      \begin{minted}[fontsize=\footnotesize]{cpp}
#include <concepts>

template<typename T>
class VecWrapper {
public:
  using Item = T;
  Item* get() { return nullptr; }
};

template<typename C>
concept Container = requires(C c) {
  typename C::Item;
  { c.get() } ->
    std::same_as<typename C::Item*>;
};
static_assert(Container<VecWrapper<int>>);

int main() {}
      \end{minted}
    \end{column}
  \end{columns}
\end{frame}

% ---------------------------------------------------------------------

\begin{frame}[fragile]{CRTP в Rust}
\begin{minted}[fontsize=\small]{rust}
trait Base<T: Base<T>> {
  fn interface(&self) {
    self.impl_method();
  }
  fn impl_method(&self);
}

struct Derived;
impl Base<Derived> for Derived {
  fn impl_method(&self) {
    println!("Derived implementation");
  }
}
\end{minted}
\end{frame}

% ---------------------------------------------------------------------

\begin{frame}[fragile]{Маркерные трейты}
\begin{minted}[fontsize=\small]{rust}
use std::rc::Rc;

// Send - threadsafity compiler marker
fn is_send<T: Send>() {}

fn main() {
  is_send::<i32>();            // OK
  is_send::<String>();         // OK
  // is_send::<Rc<i32>>();     // Error: Rc is not Send
}
\end{minted}
\end{frame}

% ---------------------------------------------------------------------

\begin{frame}[fragile]{Несколько трейт-bounds и where-клаузы}
  \begin{columns}[T, onlytextwidth]
    \setlength{\columnsep}{0pt}
    \begin{column}{0.5\textwidth}
      \begin{minted}[fontsize=\footnotesize]{rust}
use std::fmt::Debug;

fn print_if_eq<T: PartialEq + Debug>(
  a: T, b: T
) {
  if a == b { println!("{:?}", a); }
}

fn main() {
  print_if_eq(42, 42);       // 42
  print_if_eq("hi", "bye"); // (nothing)
}
      \end{minted}
    \end{column}
    \begin{column}{0.5\textwidth}
      \begin{minted}[fontsize=\footnotesize]{rust}
use std::fmt::{Debug, Display};

fn log<T, U>(val: T, label: U)
where
  T: Debug + Clone + Send,
  U: Display,
{
  println!("{}: {:?}", label, val);
}

fn main() {
  log(vec![1, 2, 3], "numbers");
  // numbers: [1, 2, 3]
}
      \end{minted}
    \end{column}
  \end{columns}
\end{frame}

% ---------------------------------------------------------------------

\begin{frame}[fragile]{\texttt{impl Trait}}
  \begin{columns}[T, onlytextwidth]
    \setlength{\columnsep}{0pt}
    \begin{column}{0.5\textwidth}
      \begin{minted}[fontsize=\footnotesize]{rust}
use std::fmt::Display;

fn print_val(val: impl Display) {
  println!("{}", val);
}

fn print_val2<T: Display>(val: T) {
  println!("{}", val);
}

fn main() {
  print_val(42);
  print_val("hello");
}
      \end{minted}
    \end{column}
    \begin{column}{0.5\textwidth}
      \begin{minted}[fontsize=\footnotesize]{rust}
fn make_adder(x: i32) -> impl Fn(i32) -> i32 {
  move |y| x + y
}

fn make_adder_dyn(x: i32) ->
    Box<dyn Fn(i32) -> i32> {
  Box::new(move |y| x + y)
}

fn main() {
  let add5 = make_adder(5);
  println!("{}", add5(3)); // 8
}
      \end{minted}
    \end{column}
  \end{columns}
\end{frame}

% ---------------------------------------------------------------------

\end{document}
